\section{Generators in a Category}

\subsection{Definition}

Let $\mathcal{C}$ be a category. An object $G \in \mathcal{C}$ is called a
\emph{generator} if for every pair of distinct morphisms
\[
f,g : X \to Y,
\]
there exists a morphism
\[
u : G \to X
\]
such that
\[
f \circ u \neq g \circ u.
\]

Equivalently, whenever
\[
f \circ u = g \circ u
\]
for all morphisms $u : G \to X$, it follows that
\[
f=g.
\]

Thus, maps from $G$ are sufficient to distinguish morphisms in the category.

\subsection{Intuition}

A generator is an object that acts as a \emph{probe}. If two morphisms differ,
then some map from the generator can detect the difference.

In many categories, morphisms from a generator correspond to the
``elements'' or ``points'' of an object.

\subsection{Examples}

\subsubsection*{1. The Category of Sets}

In the category $\mathbf{Set}$, any singleton set
\[
\{*\}
\]
is a generator.

A function
\[
u : \{*\} \to X
\]
is completely determined by the image of $*$, so there is a natural bijection
\[
\mathrm{Hom}_{\mathbf{Set}}(\{*\},X)
\cong X.
\]

If
\[
f,g : X \to Y
\]
are distinct, then there exists $x\in X$ such that
\[
f(x)\neq g(x).
\]
Define
\[
u(*)=x.
\]
Then
\[
(f\circ u)(*)=f(x)
\quad\text{and}\quad
(g\circ u)(*)=g(x),
\]
so
\[
f\circ u \neq g\circ u.
\]

Hence $\{*\}$ is a generator of $\mathbf{Set}$.

\subsubsection*{2. The Category of Groups}

In the category $\mathbf{Grp}$, the group
\[
\mathbb Z
\]
is a generator.

If
\[
f,g : G \to H
\]
are distinct group homomorphisms, then there exists
$a\in G$ such that
\[
f(a)\neq g(a).
\]

Define
\[
u:\mathbb Z \to G,
\qquad
u(n)=a^n.
\]

Then
\[
(f\circ u)(1)=f(a)
\quad\text{and}\quad
(g\circ u)(1)=g(a),
\]
so
\[
f\circ u \neq g\circ u.
\]

Therefore $\mathbb Z$ is a generator of $\mathbf{Grp}$.

\subsubsection*{3. The Category of Vector Spaces}

Let $\mathbf{Vect}_F$ denote the category of vector spaces over a field $F$.

The one-dimensional vector space
\[
F
\]
is a generator.

For every vector space $V$,
\[
\mathrm{Hom}_{\mathbf{Vect}_F}(F,V)
\cong V.
\]

Indeed, a linear map
\[
u:F\to V
\]
is determined by the vector
\[
u(1)\in V.
\]

If
\[
T,S:V\to W
\]
are distinct linear maps, there exists
$v\in V$ such that
\[
T(v)\neq S(v).
\]

Define
\[
u(a)=av.
\]

Then
\[
(T\circ u)(1)=T(v)
\quad\text{and}\quad
(S\circ u)(1)=S(v),
\]
so
\[
T\circ u \neq S\circ u.
\]

Hence $F$ is a generator of $\mathbf{Vect}_F$.

\subsubsection*{4. The Category of $R$-Modules}

Let $R\text{-}\mathbf{Mod}$ denote the category of left $R$-modules.

The module
\[
R
\]
(viewed as a left module over itself) is a generator.

For every $R$-module $M$,
\[
\mathrm{Hom}_R(R,M)
\cong M.
\]

Indeed, every module homomorphism
\[
u:R\to M
\]
is completely determined by
\[
u(1)\in M,
\]
since
\[
u(r)=r\,u(1).
\]

If
\[
f,g:M\to N
\]
are distinct module homomorphisms, then there exists
$m\in M$ such that
\[
f(m)\neq g(m).
\]

Define
\[
u:R\to M,
\qquad
u(r)=rm.
\]

Then
\[
(f\circ u)(1)=f(m)
\quad\text{and}\quad
(g\circ u)(1)=g(m),
\]
so
\[
f\circ u \neq g\circ u.
\]

Therefore $R$ is a generator of $R\text{-}\mathbf{Mod}$.

\subsubsection*{5. The Category of Topological Spaces}

In the category $\mathbf{Top}$, the one-point space
\[
\{*\}
\]
is a generator.

A continuous map
\[
u:\{*\}\to X
\]
corresponds exactly to a point of $X$.

Thus,
\[
\mathrm{Hom}_{\mathbf{Top}}(\{*\},X)
\cong X.
\]

If
\[
f,g:X\to Y
\]
are distinct continuous maps, there exists
$x\in X$ such that
\[
f(x)\neq g(x).
\]

Define
\[
u(*)=x.
\]

Then
\[
f\circ u \neq g\circ u.
\]

Hence $\{*\}$ is a generator of $\mathbf{Top}$.

\subsection{Summary Table}

\[
\begin{array}{|c|c|c|}
\hline
\text{Category} & \text{Generator} & \text{Maps from Generator Represent} \\
\hline
\mathbf{Set} & \{*\} & \text{Elements} \\
\hline
\mathbf{Grp} & \mathbb Z & \text{Group elements} \\
\hline
\mathbf{Vect}_F & F & \text{Vectors} \\
\hline
R\text{-}\mathbf{Mod} & R & \text{Module elements} \\
\hline
\mathbf{Top} & \{*\} & \text{Points} \\
\hline
\end{array}
\]

\subsection{Key Insight}

A generator allows us to study objects through morphisms from a single
distinguished object. In many categories, morphisms from the generator
correspond precisely to the ``elements'' of an object.

Thus generators convert internal structure into external morphisms, a central
theme of category theory.